% AUTORIDAD DE FUENTE: Exposé IX, Bibliografía, líneas 4743--4826;
% índices de escaneo 531--534, folios 520--523. Las capas francesas D/C
% no presentan diferencias materiales. Cinco erratas bibliográficas
% inequívocas se conservan reversiblemente mediante \sisi y se han
% comunicado al custodio raíz. Se preservan las entradas históricamente
% incompletas [13] y [24].

\bigskip
\noindent\underline{BIBLIOGRAFÍA}

\begin{enumerate}
\item[{[1]}] M. Artin, \emph{Algebraization of formal moduli I}
  (de próxima aparición).
\item[{[2]}] H. Cartan, \emph{Familles d'espaces complexes et fondements
  de la géométrie analytique}, Séminaire ENS 1960/61.
\item[{[3]}] P. Deligne, D. Mumford, \emph{The irreducibility of the space
  of curves of given genus}, Publ. Math. núm.~36 (1969).

%%%%%%%%%%  scan idx 532  |  folio 521  %%%%%%%%%%

\item[{[4]}] P. Deligne, trabajo en preparación sobre la teoría de Hodge
  (de próxima aparición en Publ. Math., 1970).
\item[{[5]}] M. Demazure, J. Giraud, M. Raynaud, \emph{Schémas
  abéliens}, Séminaire Orsay 1967/68, de próxima aparición.
\item[{[6]}] A. Grothendieck, \emph{Le Groupe de Brauer III}, en
  \emph{Dix Exposés sur la cohomologie des schémas}, North-Holland
  Publishing Co.
\item[{[7]}] A. Grothendieck, \emph{Crystals and the De Rham cohomology of
  \sisi{schémes}{schemes}} (notas de I. Coates y O. Jussila), en
  \emph{Dix exposés sur la cohomologie des schémas},
  \sisi{Noth-Holland}{North-Holland} Publishing Co.
\item[{[8]}] A. Grothendieck, \emph{Un théorème sur les homomorphismes de
  schémas abéliens}, Invent. Math. \textbf{2} (1966), pp.~59--78.
\item[{[9]}] P. A. Griffiths, \emph{Report on variation of Hodge
  structures} (de próxima aparición).
\item[{[10]}] M. Hakim, \emph{Topos annelés et schémas relatifs}, de
  próxima aparición en Grundlehren der Mathematik, Springer (1970).
\item[{[11]}] J. I. Igusa, \emph{Abstract vanishing cycle theory},
  Proceedings of the Japan Academy \textbf{34}, núm.~9 (1958),
  pp.~589--593.
\item[{[12]}] R. Kiehl, \emph{Der Endlichkeitssatz für eigentliche
  \sisi{Abbildnngen}{Abbildungen}}, Invent. Math. \textbf{2} (1967),
  pp.~191--214.
\item[{[13]}] R. Kiehl.

%%%%%%%%%%  scan idx 533  |  folio 522  %%%%%%%%%%

\item[{[14]}] S. L. Kleiman, \emph{Algebraic cycles and the Weil
  conjectures}, en \emph{Dix exposés sur la cohomologie des schémas},
  North-Holland Publishing Co.
\item[{[15]}] S. Lang, \emph{Abelian varieties}, Interscience Publishers,
  núm.~7 (1959).
\item[{[16]}] S. Lang y A. Néron, \emph{Rational points of abelian
  varieties over function fields}, Amer. J. Math. \textbf{81}, núm.~1
  (1959), pp.~95--118.
\item[{[16bis]}] U. I. Manin, \emph{Théorie des groupes formels
  commutatifs}, Uspekhi Mat. Nauk \textbf{XVIII} (1963), pp.~4--90.
\item[{[17]}] D. Mumford, \emph{Geometric invariant theory}, Ergebnisse
  der Mathematik 34, Springer (1965).
\item[{[18]}] D. Mumford, \emph{Abstract Theta functions}, Advanced
  Science Seminar in Algebraic Geometry, Bowdoin College, 1967 (notas de
  H. Pittie).
\item[{[19]}] A. Néron, \emph{Modèles minimaux des variétés abéliennes
  sur les corps locaux et globaux}, Publ. Math. 21 (1964), pp.~5--128.
\item[{[20]}] M. Raynaud, \emph{Caractéristique d'Euler--Poincaré d'un
  faisceau \sisi{et et}{et} cohomologie des variétés abéliennes}, en
  \emph{Dix exposés sur la cohomologie des schémas}, North-Holland
  Publishing Co.
\item[{[21]}] M. Raynaud, \emph{Modèles de Néron}, C. R. Acad. Sci.
  Paris \textbf{262} (febrero de 1966), pp.~413--416.
\item[{[22]}] M. Raynaud, \emph{Spécialisation du foncteur de Picard},
  C. R. Acad. Sci. \textbf{264}, pp.~941--943 y 1001--1004.
\item[{[22bis]}] M. Raynaud, \emph{Spécialisation du foncteur de
  Picard}, de próxima aparición en Publ. Math.

%%%%%%%%%%  scan idx 534  |  folio 523  %%%%%%%%%%

\item[{[23]}] M. Raynaud, \emph{Faisceaux amples sur les schémas en
  groupes et les espaces homogènes}, tesis, París, 1968 (de próxima
  aparición en Lecture Notes, Springer, 1970).
\item[{[24]}] M. Raynaud.
\item[{[25]}] M. Rosenlicht, \emph{Some basic theorems on algebraic
  groups}, Amer. J. Math. \textbf{78}, núm.~2 (1956), pp.~401--443.
\item[{[25bis]}] J.-P. Serre, \emph{Quelques propriétés des variétés
  abéliennes en caractéristique $p$}, Amer. J. Math. \textbf{80} (1958),
  pp.~715--739.
\item[{[26]}] J.-P. Serre, \emph{Corps locaux}, Act. Sci. Ind. 1296
  (1962), \sisi{Herman}{Hermann}, París.
\item[{[27]}] J.-P. Serre, \emph{Sur les corps locaux à corps résiduel
  algébriquement clos}, Bull. Soc. Math. France \textbf{89} (1961),
  pp.~105--154.
\item[{[28]}] J.-P. Serre, \emph{Abelian $\ell$-adic representations and
  elliptic curves}, Benjamin Inc. (1968).
\item[{[29]}] J.-P. Serre y J. Tate, \emph{Good reduction of abelian
  varieties}, Ann. of Math. \textbf{88}, núm.~3 (1968), pp.~492--517.
\item[{[30]}] J.-P. Serre y J. Tate, \emph{Elliptic curves and formal
  groups} (notas del Summer Institute on Algebraic Geometry, Woods Hole,
  1964).
\item[{[31]}] J.-P. Serre, \emph{Morphismes universels et variété
  d'Albanese}, Séminaire Chevalley 1958/59, Exposé 10.
\item[{[32]}] J. Tate, \emph{WC-groups over $p$-adic fields}, Séminaire
  Bourbaki núm.~156 (1957).
\item[{[33]}] J. Tate, \emph{$p$-divisible groups}, Proceedings of a
  Conference on Local Fields, Driebergen (1966), Springer.
\item[{[34]}] J. Tate, \emph{Rigid analytic spaces} (mimeografiado por
  el IHES).
\end{enumerate}
